\documentclass[12pt]{article}
\usepackage[T1]{fontenc}
\usepackage[utf8]{inputenc}
\usepackage{lmodern}
\usepackage{textcomp}
\usepackage{enumitem}
\usepackage{geometry}
\geometry{margin=1in}
\begin{document}
\begin{center}
Answer Key - Constitutional Law - Graduate Level Assessment 19 \
\small (Answers are ordered exactly as in the question paper.)
\end{center}

\section*{Multiple Choice}
\begin{enumerate}[label=\arabic*.]
\item \textbf{Answer:} A
\\ \textit{Options:} A) Protocol 14 established the Human Rights Council; B) Protocol 14 mandated a two judge formation for hearing admissibility cases; C) Protocol 14 reduced the jurisdiction of the European Court of Human Rights; D) Protocol 14 shortened the judicial term of office for the European Court of Human Rights; E) Protocol 14 introduced a requirement for unanimous decision-making in admissibility cases
\\ \textit{Note:} Protocol 14 did not establish the Human Rights Council; instead, the Council was created by the United Nations General Assembly in 2006 through a separate resolution, making the statement false.
\item \textbf{Answer:} A
\\ \textit{Options:} A) only if the prosecutor agrees.; B) immediately after being charged.; C) upon request and showing of good cause.; D) upon court order.; E) under no circumstances.
\\ \textit{Note:} The correct answer is based on the principle that access to juror information in criminal cases is often contingent upon mutual agreement between the prosecution and defense, reflecting the balance of rights and procedural fairness in the judicial process.
\item \textbf{Answer:} D
\\ \textit{Options:} A) The purchase of a commercial property building.; B) A contract for hiring a live-in nanny.; C) The purchase of stocks and bonds.; D) Crops and timber to be severed from the property next summer.; E) A common carrier delivering a new computer.
\\ \textit{Note:} The correct answer is based on the UCC's application to contracts involving the sale of goods, and since crops and timber are considered goods that can be severed from the land, they fall under the UCC's provisions.
\item \textbf{Answer:} C
\\ \textit{Options:} A) Command; B) Sociological; C) Historical; D) Interpretive; E) Analytical
\\ \textit{Note:} The Historical School of jurisprudence posits that law evolves from the customs and traditions of society, reflecting its historical context and collective values over time.
\item \textbf{Answer:} C
\\ \textit{Options:} A) Men are unable to comprehend their differences from women.'; B) Men and women have different conceptions of the feminist project.'; C) Women look to context, whereas men appeal to neutral, abstract notions of justice.'; D) 'Women and men have identical perspectives on justice.'; E) 'Feminism is about eradicating differences between men and women.'
\\ \textit{Note:} The correct answer highlights the key concept of difference feminism, which asserts that women's experiences and perspectives are shaped by their social contexts, contrasting with men's reliance on universal principles of justice, thereby emphasizing the importance of contextual understanding in feminist legal theory.
\end{enumerate}

\section*{True / False}
\begin{enumerate}[label=\arabic*.]
\setcounter{enumi}{5}
\item \textbf{Answer:} True
\\ \textit{Options:} A) True; B) False
\\ \textit{Note:} The Analytical School asserts that law is determined by the recognition and enforcement of customs by judicial bodies, thereby establishing that a custom becomes law when acknowledged by the courts.
\item \textbf{Answer:} True
\\ \textit{Options:} A) True; B) False
\\ \textit{Note:} Hart's distinction highlights that 'being obliged' refers to the external pressure to conform to a rule, while 'having an obligation' encompasses the internal acceptance of that rule, thereby clarifying how legal rules function both as social norms and as sources of legal authority in constitutional law.
\item \textbf{Answer:} True
\\ \textit{Options:} A) True; B) False
\\ \textit{Note:} Prior to the United Nations Charter, international law primarily recognized the right to use armed force solely in self-defense, as there were no comprehensive legal frameworks governing the use of force between states.
\item \textbf{Answer:} True
\\ \textit{Options:} A) True; B) False
\\ \textit{Note:} The UK Constitution is considered uncodified because it is not contained within a single, written document; instead, it is formed from various sources such as statutes, conventions, legal precedents, and works of authority, which collectively shape its constitutional framework.
\item \textbf{Answer:} False
\\ \textit{Options:} A) True; B) False
\\ \textit{Note:} Standard-setting in human rights diplomacy also encompasses the development of non-binding norms and principles that guide state behavior and promote accountability beyond mere legal enforcement.
\end{enumerate}

\section*{Long Form Response}
\begin{enumerate}[label=\arabic*.]
\setcounter{enumi}{10}
\item \textbf{Answer:} The Kadi judgment significantly reshaped the relationship between UN Security Council (UNSC) resolutions and domestic law, particularly by asserting that such resolutions must be interpreted in a manner consistent with fundamental human rights. This ruling emphasizes that domestic courts have the authority to review and potentially set aside UNSC resolutions when they conflict with established human rights norms. By prioritizing human rights standards, the Kadi case underscores the necessity for domestic legal frameworks to uphold individual rights even in the face of international obligations. Consequently, this judgment reinforces the principle that the rule of law and human rights protection must prevail, thereby promoting a more balanced approach to international law and domestic legal systems.
\\ \textit{Note:} correct because it highlights the Kadi judgment's assertion of the supremacy of fundamental human rights over UN Security Council resolutions, establishing that domestic courts can review and potentially reject such resolutions that conflict with human rights norms.
\item \textbf{Answer:} In assessing sufficient good cause or excusable neglect for a late notice of appeal, courts consider factors such as the litigant's circumstances and the reasons for the delay. For instance, if a litigant was involved in a car crash after running a red light just before the judgment was rendered and subsequently spent two months in the hospital recovering, these circumstances could be deemed as excusable neglect. The litigant's hospitalization likely impeded their ability to file an appeal within the mandated 30-day period, thereby constituting a valid reason for the delay. Courts may be sympathetic to such situations, recognizing that unforeseen medical emergencies can hinder timely legal action. Thus, in this scenario, the combination of the accident and resulting hospitalization presents a compelling case for good cause in seeking to file a late notice of appeal.
\\ \textit{Note:} correct because it illustrates how personal circumstances, such as hospitalization due to a car accident, can significantly impede a litigant's ability to meet the appeal deadline, thereby constituting excusable neglect in the eyes of the court.
\end{enumerate}

\vfill
\noindent\textit{End of Answer Key}
\end{document}